%% kapa: exam=1 round=24 year=2013 period=1 form=A subject=civil_law subject_ko=민법 seq=40 number=40
%% source: Q-Net 감정평가사 기출문제 (cst00309) artlSeq=5073802 file=extracted/1008/5073802/2013년 제24회 감정평가사 1차시험 기출문제/1교시(A형).hwp
%% conversion: kordoc 4.24.51 -> tools/gen1.py (auto); review=proofread-sonnet; flags=none
\begin{question}{40}
동산양도담보에 관한 설명으로 옳은 것은? (다툼이 있으면 판례에 의함)
\begin{choices}
  \item 점유개정의 방법으로 이중으로 양도담보를 설정한 경우 뒤의 양도담보의 채권자는 2순위 양도담보권을 취득한다.
  \item 동산양도담보에도 가등기담보 등에 관한 법률이 적용된다.
  \item 자신의 동산을 타인에게 양도담보로 제공한 자로부터 점유개정의 방법으로 그 동산을 양수한 자는 현실로 인도받지 않았더라도 그 동산을 선의취득할 수 있다.
  \item 집합물에 대해서는 양도담보를 설정할 수 없다.
  \item 신탁적 소유권이전설에 의하면 채권자는 목적 동산의 소유권을 취득한다.
\end{choices}
\end{question}
\answer{2,5}
\answernote{복수 정답(정답표 원문: 2, 5)}
